\chapter{Home Gym and Exercise (HGA)}

\section*{Definition and Overview}
HGA, used here as shorthand for home-based gym and exercise activity, refers to exercising at home using equipment such as weights, resistance bands, benches, cardio machines, and bodyweight routines. Home exercise has grown enormously in popularity and offers convenience, privacy, and flexibility. It can deliver the same health benefits as gym-based training when programmed sensibly, but it also carries risks, including injury from incorrect technique, unsafe equipment, and lack of supervision. Safe, effective home training requires appropriate equipment, a sensible plan, and attention to technique.

\section*{Epidemiology}
Participation in home-based exercise has increased substantially in recent years, with a large proportion of Australians now exercising at home, many following online programs and classes. Around half of Australian adults meet physical activity guidelines, and home exercise is helping more people to be active. Injuries from home training are common but usually minor, with muscle strains, sprains, and lower back injuries among the most frequent. Poorly secured equipment and unsupervised heavy lifting account for the more serious injuries, which are nevertheless uncommon.

\section*{Aetiology and Pathophysiology}
\begin{itemize}
    \item \textbf{Training stimulus:} resistance and aerobic exercise stress muscles, bones, and the cardiovascular system, stimulating adaptation
    \item \textbf{Muscle growth and strength:} progressive overload with adequate recovery builds muscle and strength
    \item \textbf{Metabolic benefits:} regular exercise improves blood sugar, blood pressure, cholesterol, and body composition
    \item \textbf{Injury mechanisms:} poor technique, excessive load, sudden increases in volume, fatigue, and unsafe equipment
    \item \textbf{Overtraining:} insufficient recovery causes fatigue, poor performance, and injury
    \item \textbf{Common injuries:} muscle strains, tendonitis, and lower back injury from heavy or uncontrolled lifting
\end{itemize}

\section*{Clinical Features}
\begin{itemize}
    \item Muscle soreness after new or increased training, which is normal and settles in days
    \item Sharper, persistent pain suggests injury rather than normal soreness
    \item Strains cause pain, swelling, and reduced movement in the affected muscle
    \item Lower back pain from poor lifting technique
    \item Tendon pain, such as in the shoulder or knee, with repetitive loading
    \item Falls and crush injuries from unstable equipment
    \item General health improvements: better strength, fitness, mood, and sleep
\end{itemize}

\section*{Differential Diagnosis}
\begin{itemize}
    \item Muscle soreness versus actual muscle strain or tear
    \item Tendinopathy versus bursitis and joint injury
    \item Back pain from muscle strain versus disc problems
    \item Exercise-related chest pain, which requires exclusion of cardiac causes
    \item Overtraining symptoms, which can mimic depression and chronic fatigue
    \item Any severe or persistent symptom warrants medical assessment
\end{itemize}

\section*{When to Seek Medical Care}
See a doctor for injuries that do not settle within a few days, for chest pain during exercise, or if you have health conditions that need medical clearance before starting a training program.

\textbf{Call 000 (or local emergency services) for:}
\begin{itemize}
    \item Chest pain, pressure, or severe breathlessness during or after exercise
    \item Fainting, collapse, or palpitations with dizziness
    \item A heavy weight crushing the chest or head with difficulty breathing
    \item Severe injury with a suspected fracture or spinal injury
    \item Sudden severe headache or weakness
\end{itemize}

\section*{Diagnosis and Investigations}
\begin{itemize}
    \item \textbf{Injury assessment:} clinical examination identifies most strains and sprains
    \item \textbf{Imaging:} X-ray, ultrasound, or MRI for suspected fractures, tears, or tendon injuries
    \item \textbf{Medical screening:} blood pressure, and assessment of cardiac risk before vigorous training where appropriate
    \item \textbf{Training review:} analysis of program design, technique, and progression
    \item \textbf{Equipment inspection:} checking weights, benches, and machines for safety
\end{itemize}

\section*{Management}
\begin{itemize}
    \item \textbf{Program design:} a balanced plan of resistance, cardio, and mobility training, with progressive overload
    \item \textbf{Technique first:} learn correct form, often with the help of a physiotherapist or online instruction
    \item \textbf{Injury care:} rest, ice, compression, and elevation for acute injuries, followed by graded return
    \item \textbf{Recovery:} adequate sleep, nutrition, and rest days between hard sessions
    \item \textbf{Equipment safety:} secure weights, use collars, and inspect equipment regularly
    \item \textbf{Professional input:} physiotherapists and exercise physiologists design programs for specific goals and conditions
    \item \textbf{Progress tracking:} gradual increases in load, volume, and difficulty
\end{itemize}

\section*{Complications}
\begin{itemize}
    \item Muscle and tendon injuries from poor technique or excessive load
    \item Lower back injuries from heavy, uncontrolled lifting
    \item Falls from unstable equipment or misplaced weights
    \item Overtraining syndrome with fatigue, poor performance, and mood disturbance
    \item Cardiac events in people with undiagnosed heart disease undertaking vigorous exercise
    \item Equipment-related injuries, including crush injuries
    \item Loss of motivation and cessation of exercise after injury
\end{itemize}

\section*{Prognosis and Outcomes}
Home-based exercise is safe and effective for the vast majority of people. Most minor injuries resolve within days to weeks with simple care, and serious injuries are uncommon when equipment is sound and technique is reasonable. The health benefits are substantial: regular home training improves strength, fitness, body composition, mood, and long-term health, matching gym-based training for most outcomes. Adherence is often better at home, and the main predictor of success is consistency rather than the location of training.

\section*{Prevention and Self-Care}
\begin{itemize}
    \item Start with appropriate, manageable loads and progress gradually
    \item Learn correct technique before adding heavy weights
    \item Use safe equipment: secure plates, use collars, and check benches and machines
    \item Warm up before sessions and cool down afterwards
    \item Include rest days and prioritise sleep and nutrition
    \item Keep the training area clear and well-ventilated
    \item Get medical clearance before vigorous training if you have health conditions
    \item Stop and seek advice for pain beyond normal muscle soreness
\end{itemize}

\section*{Sources and Further Reading}
The following reputable sources provide further information for healthcare students and patients seeking reliable, current advice about home-based exercise.

\begin{itemize}
    \item Healthdirect. \textit{Physical activity and exercise at home.} https://www.healthdirect.gov.au/
    \item Sports Medicine Australia. \textit{Safe home exercise and injury prevention.} https://sma.org.au/
    \item Australian Physiotherapy Association. \textit{Exercise programming and injury advice.} https://australian.physio/
    \item Exercise and Sports Science Australia. \textit{Exercise guidelines.} https://www.essa.org.au/
    \item NHS. \textit{Getting started with home exercise.} https://www.nhs.uk/live-well/
    \item World Health Organization. \textit{Physical activity guidelines.} https://www.who.int/
\end{itemize}
